
\documentclass{article}
\usepackage{colortbl}
\usepackage{makecell}
\usepackage{multirow}
\usepackage{supertabular}

\begin{document}

\newcounter{utterance}

\centering \large Interaction Transcript for game `wordle', experiment `medium\_frequency\_words\_no\_clue\_no\_critic', episode 4 with gpt{-}5.4{-}azure.
\vspace{24pt}

{ \footnotesize  \setcounter{utterance}{1}
\setlength{\tabcolsep}{0pt}
\begin{supertabular}{c@{$\;$}|p{.15\linewidth}@{}p{.15\linewidth}p{.15\linewidth}p{.15\linewidth}p{.15\linewidth}p{.15\linewidth}}
    \# & \multicolumn{2}{c}{Player} && \multicolumn{2}{c}{Game Master} \\
    \hline

    \theutterance \stepcounter{utterance}  
    & & & \multicolumn{4}{p{0.6\linewidth}}{
        \cellcolor[rgb]{0.9,0.9,0.9}{
            \makecell[{{p{\linewidth}}}]{
                \texttt{\tiny{[P1$\langle$GM]}}
                \texttt{Tu esi kalbos burtininkas, kuriam patinka spėti žodžius, laikantis pateiktų taisyklių.} \\
\\ 
\texttt{Sveiki atvykę į Wordle! Turite šešis bandymus atspėti tikslinį žodį – galiojantį penkių mažųjų raidžių anglų kalbos žodį. Prašome naudoti žymas „paaiškinimas:" ir „spėjimas:", kad pateiktumėte glaustą paaiškinimą kiekvienam spėjimui.} \\
\\ 
\texttt{Pavyzdžiui, jei jūsų spėjimas yra „įmonė", jūsų atsakymas turėtų būti} \\
\texttt{paaiškinimas: tai yra dažnas penkių raidžių anglų kalbos žodis, ir savo spėjimą pradedu šiuo žodžiu.} \\
\texttt{spėjimas: įmonė} \\
\\ 
\texttt{Po kiekvieno spėjimo jūsų atsakymas bus patikrintas ir gausite grįžtamąjį ryšį, nurodantį, kurios raidės yra teisingos (žalios), kurios raidės teisingos, bet netinkamoje pozicijoje (geltonos), ir kurios raidės neteisingos (raudonos). Šis grįžtamasis ryšys gali būti naudingas sprendžiant, kurias raides įtraukti arba atmesti kitame spėjime.} \\
\\ 
\texttt{Pavyzdžiui, grįžtamasis ryšys žodžiui „įmonė" galėtų būti:} \\
\texttt{spėjimo\_atsiliepimas: į<geltona> m<geltona> o<žalia> n<geltona> ė<raudona>} \\
\\ 
\texttt{Paaiškinime turi būti detalių apie tai, kaip spėjimo\_atsiliepimas naudojamas naujam spėjimui suformuoti.} \\
\\ 
\texttt{Negeneruokite jokio kito teksto.} \\
\\ 
\texttt{Pradėkime nuo jūsų pirmojo spėjimo.} \\
            }
        }
    }
     \\ \\

    \theutterance \stepcounter{utterance}  
    & \multicolumn{4}{p{0.6\linewidth}}{
        \cellcolor[rgb]{0.561,0.737,0.561}{
            \makecell[{{p{\linewidth}}}]{
                \texttt{\tiny{[P1$\rangle$GM]}}
                \texttt{paaiškinimas: pradedu nuo dažno penkių raidžių anglų kalbos žodžio su dažnomis raidėmis, kad greitai patikrinčiau balses ir paplitusius priebalsius.} \\
\texttt{spėjimas: stare} \\
            }
        }
    }
    & & \\ \\

    \theutterance \stepcounter{utterance}  
    & & & \multicolumn{2}{p{0.3\linewidth}}{
        \cellcolor[rgb]{0.9,0.9,0.9}{
            \makecell[{{p{\linewidth}}}]{
                \texttt{\tiny{[GM$|$GM]}}
                \texttt{attempts = 1} \\
\texttt{target\_word = ugnis} \\
\texttt{guess = stare} \\
\texttt{guess\_feedback = s<yellow> t<red> a<red> r<red> e<red>} \\
            }
        }
    }
    & & \\ \\

    \theutterance \stepcounter{utterance}  
    & & & \multicolumn{4}{p{0.6\linewidth}}{
        \cellcolor[rgb]{0.9,0.9,0.9}{
            \makecell[{{p{\linewidth}}}]{
                \texttt{\tiny{[P1$\langle$GM]}}
                \texttt{spėjimo\_atsiliepimas: s<geltona> t<raudona> a<raudona> r<raudona> e<raudona>} \\
\\ 
\texttt{Atsakymą pateikite tik šiuo formatu.} \\
\texttt{paaiškinimas: detalės} \\
\texttt{spėjimas: žodis} \\
            }
        }
    }
     \\ \\

    \theutterance \stepcounter{utterance}  
    & \multicolumn{4}{p{0.6\linewidth}}{
        \cellcolor[rgb]{0.561,0.737,0.561}{
            \makecell[{{p{\linewidth}}}]{
                \texttt{\tiny{[P1$\rangle$GM]}}
                \texttt{paaiškinimas: iš „stare“ žinau, kad raidė s yra žodyje, bet negali būti pirmoje pozicijoje, o t, a, r ir e reikia atmesti. Todėl renkuosi žodį su s kitoje vietoje ir keturiomis naujomis dažnomis raidėmis, kad surinkčiau daugiau informacijos.} \\
\texttt{spėjimas: music} \\
            }
        }
    }
    & & \\ \\

    \theutterance \stepcounter{utterance}  
    & & & \multicolumn{2}{p{0.3\linewidth}}{
        \cellcolor[rgb]{0.9,0.9,0.9}{
            \makecell[{{p{\linewidth}}}]{
                \texttt{\tiny{[GM$|$GM]}}
                \texttt{Error: The guessed word is not a valid word for this game.} \\
            }
        }
    }
    & & \\ \\

    \theutterance \stepcounter{utterance}  
    & & & \multicolumn{4}{p{0.6\linewidth}}{
        \cellcolor[rgb]{0.9,0.9,0.9}{
            \makecell[{{p{\linewidth}}}]{
                \texttt{\tiny{[P1$\langle$GM]}}
                \texttt{Jūsų spėjimas nėra tinkamas žodis šiam žaidimui. Pabandykite dar kartą.} \\
\\ 
\texttt{Atsakymą pateikite tik šiuo formatu.} \\
\texttt{paaiškinimas: detalės} \\
\texttt{spėjimas: žodis} \\
            }
        }
    }
     \\ \\

    \theutterance \stepcounter{utterance}  
    & \multicolumn{4}{p{0.6\linewidth}}{
        \cellcolor[rgb]{0.561,0.737,0.561}{
            \makecell[{{p{\linewidth}}}]{
                \texttt{\tiny{[P1$\rangle$GM]}}
                \texttt{paaiškinimas: iš „stare“ vis dar naudoju tik tai, kad s yra žodyje, bet ne pirmoje vietoje, o t, a, r ir e netinka. Kadangi ankstesnis bandymas nebuvo priimtas, dabar renkuosi labai įprastą penkių raidžių anglų kalbos žodį, kuriame s yra penktoje pozicijoje ir tikrinamos naujos raidės l, i, o ir n.} \\
\texttt{spėjimas: lions} \\
            }
        }
    }
    & & \\ \\

    \theutterance \stepcounter{utterance}  
    & & & \multicolumn{2}{p{0.3\linewidth}}{
        \cellcolor[rgb]{0.9,0.9,0.9}{
            \makecell[{{p{\linewidth}}}]{
                \texttt{\tiny{[GM$|$GM]}}
                \texttt{Error: The guessed word is not a valid word for this game.} \\
            }
        }
    }
    & & \\ \\

    \theutterance \stepcounter{utterance}  
    & & & \multicolumn{4}{p{0.6\linewidth}}{
        \cellcolor[rgb]{0.9,0.9,0.9}{
            \makecell[{{p{\linewidth}}}]{
                \texttt{\tiny{[P1$\langle$GM]}}
                \texttt{Jūsų spėjimas nėra tinkamas žodis šiam žaidimui. Pabandykite dar kartą.} \\
\\ 
\texttt{Atsakymą pateikite tik šiuo formatu.} \\
\texttt{paaiškinimas: detalės} \\
\texttt{spėjimas: žodis} \\
            }
        }
    }
     \\ \\

    \theutterance \stepcounter{utterance}  
    & \multicolumn{4}{p{0.6\linewidth}}{
        \cellcolor[rgb]{0.561,0.737,0.561}{
            \makecell[{{p{\linewidth}}}]{
                \texttt{\tiny{[P1$\rangle$GM]}}
                \texttt{paaiškinimas: iš „stare“ naudoju tai, kad s yra žodyje, bet ne pirmoje pozicijoje, o t, a, r ir e atmetu. Todėl renkuosi dažną penkių raidžių anglų kalbos žodį, kuriame s yra ketvirtoje vietoje, o naujos raidės n, o, i ir y padeda patikrinti daugiau galimų variantų.} \\
\texttt{spėjimas: noisy} \\
            }
        }
    }
    & & \\ \\

    \theutterance \stepcounter{utterance}  
    & & & \multicolumn{2}{p{0.3\linewidth}}{
        \cellcolor[rgb]{0.9,0.9,0.9}{
            \makecell[{{p{\linewidth}}}]{
                \texttt{\tiny{[GM$|$GM]}}
                \texttt{Error: The guessed word is not a valid word for this game.} \\
            }
        }
    }
    & & \\ \\

    \theutterance \stepcounter{utterance}  
    & & & \multicolumn{4}{p{0.6\linewidth}}{
        \cellcolor[rgb]{0.9,0.9,0.9}{
            \makecell[{{p{\linewidth}}}]{
                \texttt{\tiny{[P1$\langle$GM]}}
                \texttt{Jūsų spėjimas nėra tinkamas žodis šiam žaidimui. Pabandykite dar kartą.} \\
\\ 
\texttt{Atsakymą pateikite tik šiuo formatu.} \\
\texttt{paaiškinimas: detalės} \\
\texttt{spėjimas: žodis} \\
            }
        }
    }
     \\ \\

    \theutterance \stepcounter{utterance}  
    & \multicolumn{4}{p{0.6\linewidth}}{
        \cellcolor[rgb]{0.561,0.737,0.561}{
            \makecell[{{p{\linewidth}}}]{
                \texttt{\tiny{[P1$\rangle$GM]}}
                \texttt{paaiškinimas: iš „stare“ naudoju tai, kad s yra žodyje, bet ne pirmoje pozicijoje, o t, a, r ir e netinka. Todėl renkuosi dažną penkių raidžių anglų kalbos žodį „husky“, kuriame s yra trečioje vietoje ir tikrinamos naujos raidės h, u, k ir y.} \\
\texttt{spėjimas: husky} \\
            }
        }
    }
    & & \\ \\

    \theutterance \stepcounter{utterance}  
    & & & \multicolumn{2}{p{0.3\linewidth}}{
        \cellcolor[rgb]{0.9,0.9,0.9}{
            \makecell[{{p{\linewidth}}}]{
                \texttt{\tiny{[GM$|$GM]}}
                \texttt{Error: The guessed word is not a valid word for this game.} \\
            }
        }
    }
    & & \\ \\

    \theutterance \stepcounter{utterance}  
    & & & \multicolumn{4}{p{0.6\linewidth}}{
        \cellcolor[rgb]{0.9,0.9,0.9}{
            \makecell[{{p{\linewidth}}}]{
                \texttt{\tiny{[P1$\langle$GM]}}
                \texttt{Jūsų spėjimas nėra tinkamas žodis šiam žaidimui. Pabandykite dar kartą.} \\
\\ 
\texttt{Atsakymą pateikite tik šiuo formatu.} \\
\texttt{paaiškinimas: detalės} \\
\texttt{spėjimas: žodis} \\
            }
        }
    }
     \\ \\

    \theutterance \stepcounter{utterance}  
    & \multicolumn{4}{p{0.6\linewidth}}{
        \cellcolor[rgb]{0.561,0.737,0.561}{
            \makecell[{{p{\linewidth}}}]{
                \texttt{\tiny{[P1$\rangle$GM]}}
                \texttt{paaiškinimas: iš „stare“ žinau, kad s yra žodyje, bet ne pirmoje pozicijoje, o t, a, r ir e reikia atmesti; ankstesni nepriimti bandymai nesuteikė naujos informacijos. Todėl renkuosi dažną penkių raidžių anglų kalbos žodį, kuriame s yra penktoje vietoje, o naujos raidės m, i, n ir u padeda toliau siaurinti variantus.} \\
\texttt{spėjimas: minus} \\
            }
        }
    }
    & & \\ \\

    \theutterance \stepcounter{utterance}  
    & & & \multicolumn{2}{p{0.3\linewidth}}{
        \cellcolor[rgb]{0.9,0.9,0.9}{
            \makecell[{{p{\linewidth}}}]{
                \texttt{\tiny{[GM$|$GM]}}
                \texttt{attempts = 2} \\
\texttt{target\_word = ugnis} \\
\texttt{guess = minus} \\
\texttt{guess\_feedback = m<red> i<yellow> n<green> u<yellow> s<green>} \\
            }
        }
    }
    & & \\ \\

    \theutterance \stepcounter{utterance}  
    & & & \multicolumn{4}{p{0.6\linewidth}}{
        \cellcolor[rgb]{0.9,0.9,0.9}{
            \makecell[{{p{\linewidth}}}]{
                \texttt{\tiny{[P1$\langle$GM]}}
                \texttt{spėjimo\_atsiliepimas: m<raudona> i<geltona> n<žalia> u<geltona> s<žalia>} \\
\\ 
\texttt{Atsakymą pateikite tik šiuo formatu.} \\
\texttt{paaiškinimas: detalės} \\
\texttt{spėjimas: žodis} \\
            }
        }
    }
     \\ \\

    \theutterance \stepcounter{utterance}  
    & \multicolumn{4}{p{0.6\linewidth}}{
        \cellcolor[rgb]{0.561,0.737,0.561}{
            \makecell[{{p{\linewidth}}}]{
                \texttt{\tiny{[P1$\rangle$GM]}}
                \texttt{paaiškinimas: iš „minus“ žinau, kad n yra trečioje vietoje, s yra penktoje vietoje, i yra žodyje, bet ne antroje vietoje, o u yra žodyje, bet ne ketvirtoje vietoje; taip pat m atmetu. Kartu su ankstesniu atsiliepimu atmetu t, a, r ir e, todėl renkuosi žodį, kuris tiksliai atitinka šabloną \_ u n i s ir išlaiko s ne pirmoje pozicijoje.} \\
\texttt{spėjimas: funis} \\
            }
        }
    }
    & & \\ \\

    \theutterance \stepcounter{utterance}  
    & & & \multicolumn{2}{p{0.3\linewidth}}{
        \cellcolor[rgb]{0.9,0.9,0.9}{
            \makecell[{{p{\linewidth}}}]{
                \texttt{\tiny{[GM$|$GM]}}
                \texttt{Error: The guessed word is not a valid word for this game.} \\
            }
        }
    }
    & & \\ \\

    \theutterance \stepcounter{utterance}  
    & & & \multicolumn{4}{p{0.6\linewidth}}{
        \cellcolor[rgb]{0.9,0.9,0.9}{
            \makecell[{{p{\linewidth}}}]{
                \texttt{\tiny{[P1$\langle$GM]}}
                \texttt{Jūsų spėjimas nėra tinkamas žodis šiam žaidimui. Pabandykite dar kartą.} \\
\\ 
\texttt{Atsakymą pateikite tik šiuo formatu.} \\
\texttt{paaiškinimas: detalės} \\
\texttt{spėjimas: žodis} \\
            }
        }
    }
     \\ \\

    \theutterance \stepcounter{utterance}  
    & \multicolumn{4}{p{0.6\linewidth}}{
        \cellcolor[rgb]{0.561,0.737,0.561}{
            \makecell[{{p{\linewidth}}}]{
                \texttt{\tiny{[P1$\rangle$GM]}}
                \texttt{paaiškinimas: iš „minus“ gaunu tvirtą šabloną: n yra trečioje vietoje, s penktoje, i turi būti ne antroje vietoje, o u ne ketvirtoje; taip pat atmetu m, o iš pirmojo atsiliepimo dar atmetu t, a, r ir e. Todėl renkuosi žodį, atitinkantį modelį \_ u n i s ir išlaikantį visas žinomas sąlygas.} \\
\texttt{spėjimas: lunis} \\
            }
        }
    }
    & & \\ \\

    \theutterance \stepcounter{utterance}  
    & & & \multicolumn{2}{p{0.3\linewidth}}{
        \cellcolor[rgb]{0.9,0.9,0.9}{
            \makecell[{{p{\linewidth}}}]{
                \texttt{\tiny{[GM$|$GM]}}
                \texttt{Error: The guessed word is not a valid word for this game.} \\
            }
        }
    }
    & & \\ \\

    \theutterance \stepcounter{utterance}  
    & & & \multicolumn{4}{p{0.6\linewidth}}{
        \cellcolor[rgb]{0.9,0.9,0.9}{
            \makecell[{{p{\linewidth}}}]{
                \texttt{\tiny{[P1$\langle$GM]}}
                \texttt{Jūsų spėjimas nėra tinkamas žodis šiam žaidimui. Pabandykite dar kartą.} \\
\\ 
\texttt{Atsakymą pateikite tik šiuo formatu.} \\
\texttt{paaiškinimas: detalės} \\
\texttt{spėjimas: žodis} \\
            }
        }
    }
     \\ \\

    \theutterance \stepcounter{utterance}  
    & \multicolumn{4}{p{0.6\linewidth}}{
        \cellcolor[rgb]{0.561,0.737,0.561}{
            \makecell[{{p{\linewidth}}}]{
                \texttt{\tiny{[P1$\rangle$GM]}}
                \texttt{paaiškinimas: iš „minus“ lieka labai siauras modelis: n yra trečioje vietoje, s yra penktoje, i turi būti ne antroje vietoje, o u turi būti ne ketvirtoje; taip pat atmetu m, t, a, r ir e. Todėl renkuosi žodį, atitinkantį šabloną \_unis, kad išlaikyčiau visas patvirtintas pozicijas ir patikrinčiau naują pirmąją raidę.} \\
\texttt{spėjimas: gunis} \\
            }
        }
    }
    & & \\ \\

    \theutterance \stepcounter{utterance}  
    & & & \multicolumn{2}{p{0.3\linewidth}}{
        \cellcolor[rgb]{0.9,0.9,0.9}{
            \makecell[{{p{\linewidth}}}]{
                \texttt{\tiny{[GM$|$GM]}}
                \texttt{Error: The guessed word is not a valid word for this game.} \\
            }
        }
    }
    & & \\ \\

    \theutterance \stepcounter{utterance}  
    & & & \multicolumn{4}{p{0.6\linewidth}}{
        \cellcolor[rgb]{0.9,0.9,0.9}{
            \makecell[{{p{\linewidth}}}]{
                \texttt{\tiny{[P1$\langle$GM]}}
                \texttt{Jūsų spėjimas nėra tinkamas žodis šiam žaidimui. Pabandykite dar kartą.} \\
\\ 
\texttt{Atsakymą pateikite tik šiuo formatu.} \\
\texttt{paaiškinimas: detalės} \\
\texttt{spėjimas: žodis} \\
            }
        }
    }
     \\ \\

    \theutterance \stepcounter{utterance}  
    & \multicolumn{4}{p{0.6\linewidth}}{
        \cellcolor[rgb]{0.561,0.737,0.561}{
            \makecell[{{p{\linewidth}}}]{
                \texttt{\tiny{[P1$\rangle$GM]}}
                \texttt{paaiškinimas: iš „minus“ gaunu aiškų šabloną: n yra trečioje vietoje, s yra penktoje, i yra žodyje, bet ne antroje vietoje, todėl patogiausia dėti ją į ketvirtą, o u yra žodyje, bet ne ketvirtoje vietoje, todėl ją dedu į pirmą. Taip pat atmetu m, t, a, r ir e. Žodis „ugnis“ tiksliai atitinka modelį u \_ n i s ir visas gautas sąlygas.} \\
\texttt{spėjimas: ugnis} \\
            }
        }
    }
    & & \\ \\

    \theutterance \stepcounter{utterance}  
    & & & \multicolumn{2}{p{0.3\linewidth}}{
        \cellcolor[rgb]{0.9,0.9,0.9}{
            \makecell[{{p{\linewidth}}}]{
                \texttt{\tiny{[GM$|$GM]}}
                \texttt{attempts = 3} \\
\texttt{target\_word = ugnis} \\
\texttt{guess = ugnis} \\
\texttt{guess\_feedback = u<green> g<green> n<green> i<green> s<green>} \\
            }
        }
    }
    & & \\ \\

    \theutterance \stepcounter{utterance}  
    & & & \multicolumn{2}{p{0.3\linewidth}}{
        \cellcolor[rgb]{0.9,0.9,0.9}{
            \makecell[{{p{\linewidth}}}]{
                \texttt{\tiny{[GM$|$GM]}}
                \texttt{game\_result = WIN} \\
            }
        }
    }
    & & \\ \\

\end{supertabular}
}

\end{document}
